\documentclass[11pt,a4paper]{article}
\usepackage[utf8]{inputenc}
\usepackage[T1]{fontenc}
\usepackage{lmodern}
\usepackage{amsmath,amssymb}
\usepackage[colorlinks=true,linkcolor=blue,citecolor=blue,urlcolor=blue]{hyperref}
\usepackage{geometry}
\usepackage[protrusion=true,expansion=false]{microtype}
\geometry{margin=2.8cm}

\title{\textbf{Related Work}\\[0.8ex]
\large Principia: An Interactive Instrument\\
for the Three-Body Initial-Condition Manifold}
\author{Malachy Doherty}
\date{March 2026, revised October 2026}

\begin{document}
\maketitle
\tableofcontents
\vspace{1em}

% =========================================================
\section{Introduction}
% =========================================================

Principia is an instrument for exploring the initial-condition (IC) manifold of the planar
gravitational three-body problem interactively.  It renders a zoomable two-dimensional colour
map whose every pixel is one symmetry-reduced IC, integrated forward in time, and whose colour
encodes the outcome of that trajectory or a per-trajectory diagnostic.  It is a native
application first: one Rust engine whose integration kernel is compiled to single precision for
the GPU and to double precision for the CPU, with a desktop interface; the first web version
compiles that whole native application to run in the browser.  The Burrau--Pythagorean research
programme uses Principia as its computational platform, for a systematic survey of the infinite
family of configurations indexed by primitive Pythagorean triples and of the mass and momentum
space around each member of that family.

This document surveys the research traditions on which Principia draws---IC-manifold mapping,
shape-space geometry, symbolic dynamics of triple encounters, final-motion classification, chaos
diagnostics, statistical theory of outcomes, regularised and symplectic integration, numerical
reliability, and GPU computation---and identifies the gaps in the existing literature that
motivate the present work.

Several constructs are proposed as part of the present work rather than as established
literature: a topological encounter word computed for every initial condition rather than only
for periodic orbits (\S\ref{sec:symbolic}); a ``projective microscope'' that rotates the 2D
rendering plane within the 8D manifold (\S\ref{sec:number_theory}); and the use of the
Pythagorean-triple family as a structured indexing scheme for three-body IC space
(\S\ref{sec:number_theory}).  These are labelled as proposed contributions wherever they appear.

A note on the object of study is useful at the outset.  The planar three-body problem has a
12-dimensional phase space.  Removing the centre of mass (position and velocity) leaves 8;
letting the three masses vary on the simplex $\sum m_i = 1$ adds 2; quotienting by rotation and
by the scaling symmetry of Newtonian gravity removes 2 more.  The space of distinct initial
conditions is therefore 8-dimensional: two configuration (shape) coordinates, four momentum
coordinates, and two mass-ratio coordinates.  Total energy and angular momentum are not fixed:
they are functions on this space, derived from each IC rather than imposed as constraints.  All
prior pixel-scale IC surveys are restricted to a 2-dimensional slice of this space (equal masses,
zero or near-zero velocities).  Principia is motivated by the need to navigate \emph{families} of
such 2D slices inside the full 8-dimensional manifold, making the structure of that
manifold---not merely one or two cross-sections of it---the primary object of study.

% =========================================================
\section{The Three-Body Problem: Fundamentals and History}
% =========================================================

The three-body problem---determining the motion of three mutually gravitating point masses---was
formulated in its modern form by Newton in the \emph{Principia Mathematica} of
1687~\cite{newton1687}.  Despite the simplicity of the governing equations, only exceptional
families of closed-form solutions are known: Euler's collinear configurations (1763) and
Lagrange's equilateral configurations (1772).  No general analytic solution is known.

The non-existence of further algebraic first integrals beyond the classical ones was established
by Bruns~\cite{bruns1887}, and Poincar\'{e}'s \emph{M\'{e}thodes Nouvelles de la M\'{e}canique
C\'{e}leste}~\cite{poincare1892} showed that the problem is generically non-integrable: small
perturbations of periodic solutions lead to homoclinic tangles and chaotic behaviour.  Fractal
basin boundaries in the IC space are what this predicts.

A convergent power-series solution was given by Sundman~\cite{sundman1912}, but it converges too
slowly for practical computation.  The problem must therefore be treated numerically.
Comprehensive introductions from the celestial-mechanics perspective are given by
Szebehely~\cite{szebehely1967book}, Valtonen et al.~\cite{valtonen2016}, and Heggie \&
Hut~\cite{heggie2003}.

Generic bound three-body systems almost always dissolve into a binary and an escaping single
body~\cite{heggie2003}; understanding how the outcome depends on initial conditions, and charting
the fractal structure of that dependence across the full IC manifold, is the central problem the
work described here addresses.

% =========================================================
\section{The Pythagorean Three-Body Problem}
% =========================================================

The configuration now called the \emph{Pythagorean three-body problem} was introduced by
Burrau in 1913~\cite{burrau1913}: three bodies with masses $m_1=3$, $m_2=4$, $m_3=5$ placed
at rest at the vertices of a 3--4--5 right triangle, with each mass equal to the length of the
opposite side.  Szebehely \& Peters~\cite{szebehely1967} gave the first complete numerical
solution using regularisation: after a protracted series of close triple approaches, the
lightest body escapes, leaving a bound binary of the two heavier bodies.  This trajectory has
since been a canonical benchmark for three-body codes and a standard demonstration of
sensitivity to initial conditions.

Aarseth, Anosova, Orlov \& Szebehely~\cite{aarseth1994} investigated the chaoticity structure
near the standard Pythagorean IC.  They displaced one body at a time to 100 points on each of six
circles of radius $5 \times 10^{-8}$ to $5 \times 10^{-3}$ about its standard position, recorded
the escape angle and the identity of the escaper, and used time-reversed integrations to set
aside non-predictable runs.  A small neighbourhood of the standard position of each body is a
regular region, in which the final outcome varies continuously with the initial position; other
regular regions exist, and outside them the outcome changes abruptly.  Escape follows a close
triple approach, and close triple approaches often made the motion non-predictable even with
regularisation---a first map of the IC-space structure near the Pythagorean point, and an early
statement of the reliability problem taken up in \S\ref{sec:reliability}.  The Burrau problem
remains a test case: Boekholt, Moerman \& Portegies Zwart~\cite{boekholt2021} solve it to
convergence in its Newtonian limit as the control for a relativistic study, and Chitan, Myll\"{a}ri
\& Haque~\cite{chitan2022} integrate all sixteen Pythagorean triangles with hypotenuse below 100,
each mass equal to the opposite side, mainly to measure post-Newtonian effects.

The (3,4,5) configuration is the simplest member of an infinite family indexed by primitive
Pythagorean triples: every primitive triple $(a,b,c)$ with $a^2+b^2=c^2$ yields a distinct
right-triangle geometry, and the Burrau mass convention (mass equal to the opposite side,
normalised) yields a distinct three-body configuration with rationally related mass ratios.
The number-theoretic structure of this family and its implications for the research programme
are discussed in Section~\ref{sec:number_theory}.

% =========================================================
\section{Mapping the Initial-Condition Manifold}
% =========================================================

\subsection{The Agekyan--Anosova map}

The earliest systematic attempt to sample the IC space of the general three-body problem is due
to Agekyan \& Anosova~\cite{agekian1967}.  For equal masses released (almost) from rest, the
configuration is determined, up to scale and rotation, by the shape of the triangle.  They fixed
two bodies at $(\mp\tfrac12, 0)$ and let the third range over a region $D$ bounded by the
positive coordinate axes and the unit circle centred on the first body,
$(\xi+\tfrac12)^2+\eta^2 \le 1$ with $\xi,\eta \ge 0$, which contains every configuration once.
Their 100 trajectories, with ICs drawn at random in $D$, were given small equal and opposite
initial velocities rather than none, to avoid a collision at the first close approach; they
tabulated each IC with its energy, angular momentum, disintegration time and the perimeter of the
closest triple approach.  Every completed computation ended in disintegration after a close
triple approach.  A companion paper~\cite{agekian1968b} continued the study and reprints the initial
conditions.  This work introduced the paradigm of treating IC space as a
two-dimensional region to be surveyed by outcome, and region $D$ underpins all subsequent
IC-manifold mapping in this tradition.  Li \& Liao~\cite{li2018} restate region $D$ at exact
rest and use it to search for periodic free-fall orbits.

A key limitation of the Agekyan--Anosova map is that it fixes six of the eight dimensions: equal
masses and (near-)zero initial velocities.  This surface is a codimension-6 submanifold of the
8-dimensional symmetry-reduced IC space.  The dynamics in the remaining six dimensions
(independent mass variation and non-zero momenta) are not represented.

\subsection{The Lehto et al.\ pixel map and the Maasai shield}

Lehto, Kotiranta, Valtonen, Heinäm\"{a}ki, Mikkola \& Chernin~\cite{lehto2008} carried out a
pixel-by-pixel computation over the Agekyan--Anosova domain using approximately 200,000
simulations, colour-coding each IC by the lifetime of the triple system before dissolution.  The
resulting image resembles a Maasai warrior's shield: broad bands of short-lived systems flanked
by intricate swirling structures of longer-lived systems.  Zooming in reveals the same band
structure repeating at every scale, consistent with fractal basin boundaries; Lehto et al.\ report
a Hausdorff dimension of 1.5 for their map and find 28.6\% of positions locally
homogeneous at its finest scale, and they map time-reversibility over the same domain.

This study is the closest existing precedent to Principia's rendering task.  It is, however,
subject to the same restrictions as the Agekyan--Anosova map: equal masses, zero velocities, and
a single diagnostic (lifetime), computed in batch with no provision for interactive navigation or
zoom.  For Principia, the Lehto et al.\ study both validates the rendering paradigm and defines
the narrow 2D subspace within which all prior high-resolution work has been done.

\subsection{Converged maps and the structure of the free-fall plane}

Boekholt, Portegies Zwart \& Valtonen~\cite{boekholt2020} integrated a uniform grid of 1212
initial conditions over the Agekyan--Anosova domain with the arbitrary-precision code Brutus, at
accuracies from $10^{-6}$ down to $10^{-70}$, to $t = 10^4$.  The same initial condition can live
for very different times at the two extremes while the two lifetime maps remain statistically
indistinguishable, and the fraction of solutions that cannot be reversed falls as a power law
with the numerical accuracy: a map of this domain whose individual pixels have a known accuracy.

Tanikawa and collaborators studied the same plane through its collisions.  The initial
conditions that lead to a binary collision form curves, and these curves bound the regions of
the plane that share a symbol sequence~\cite{tanikawa1995,tanikawa2008,tanikawa2015}; Tanikawa
\& Mikkola~\cite{tanikawa2021} confirm numerically that triple-collision orbits exist inside the
domain, at points where binary-collision curves meet, and tabulate them.  Mikkola \&
Tanikawa~\cite{mikkola2007} relate the lifetime of each free-fall system to its Lyapunov time,
which ties a lifetime map and a Lyapunov map of the same plane together.

Parischewsky, Trani \& Leigh~\cite{parischewsky2023} map planar free fall near the equilateral
(Lagrange) configuration, from the triple-collision singularity out to the hierarchical limit,
with outcome maps in the space of side ratios that separate regular from ergodic runs and compare
both with the statistical escape theory.

\subsection{Beyond the zero-velocity, equal-mass surface}

Published maps off this surface are few and narrow.  Tanikawa \&
Mikkola~\cite{tanikawa2008} extend their first-symbol maps of the free-fall plane to non-zero
angular momentum.  \v{S}uvakov \& Dmitra\v{s}inovi\'{c}~\cite{suvakov2013} search a
two-dimensional slice of initial velocities, at fixed collinear positions and equal masses, for
periodic orbits.  Li, Li \& Liao~\cite{lili2020} continue one stable periodic-orbit family across
a plane of two mass ratios and map its period and stability there---the closest published
analogue to a mass chart, though for one orbit family rather than for outcomes.  Anosova~\cite{anosova1969}
gave the first outcome statistics for unequal masses in region $D$, and Valtonen, Myll\"{a}ri,
Orlov \& Rubinov~\cite{valtonen2005} integrated 120,000 rotating equal-mass systems from random
initial conditions, but these are statistics over random samples, not maps.

\subsection{Trani et al.\ and isles of regularity}

Trani, Leigh, Boekholt \& Portegies Zwart~\cite{trani2024} performed a large-scale IC scan with
the \emph{Tsunami} regularised N-body code, running $\mathcal{O}(10^6)$ scattering simulations
with unequal masses over a 2D parameterisation that varies the phase of a pre-existing binary and
the inclination of the incoming third body.  They found that regular (non-chaotic) trajectories occupy a substantial
fraction of the scanned IC space, forming connected ``isles of regularity'' in the chaotic sea,
and they separate physical chaos from numerical chaos in doing so.  This confirms the picture of
a fractal IC-space boundary interspersed with structured regular regions, and demonstrates that
modern hardware can sustain IC surveys at million-trajectory scale.  Principia's own prototype
finds both regimes within a single slice: regions where neighbouring ICs are the same trajectory
shifted in orbital phase, and regions where they decorrelate.

The scattering parameterisation of Trani et al.\ is adapted to hierarchical encounters (one body
arriving from far away to interact with a pre-existing binary) and does not cover the
non-hierarchical regime in which all three bodies start with comparable separations; its
inclination axis also takes the map out of the plane, so only one row of it lies in the planar
problem.  It is also
static: results are computed offline and cannot be explored interactively.

\subsection{Fractal basin boundaries}
\label{sec:fractal}

The theoretical framework for fractal basin boundaries was established by McDonald, Grebogi, Ott
\& Yorke~\cite{mcdonald1985}, who showed that the boundary between basins of competing attractors
is generically fractal when the underlying dynamics are chaotic.  Aguirre, Vallejo \&
Sanju\'{a}n~\cite{aguirre2001} showed that in open Hamiltonian systems with three or more escape
channels the basin boundaries often satisfy the stronger \emph{Wada} property: every boundary
point is a limit point of three or more basins simultaneously.  For the three-body problem the
relevant escape channels are which body escapes and in which direction; whether its basin
boundaries are Wada is a question the IC maps can address directly.

Two tools are standard for quantifying such boundaries.  The \emph{uncertainty exponent}
(Grebogi, McDonald, Ott \& Yorke~\cite{grebogi1983}; see also~\cite{mcdonald1985}) measures how
the fraction of initial conditions whose outcome flips under an $\varepsilon$-perturbation scales
with $\varepsilon$, yielding the uncertainty dimension of the boundary; \emph{basin entropy}
(Daza, Wagemakers, Georgeot, Gu\'{e}ry-Odelin \& Sanju\'{a}n~\cite{daza2016}) measures the Shannon
uncertainty of outcomes over a grid and, in its boundary form, gives a test for the Wada property.
These are the natural instruments for the IC-space boundaries mapped here, and they carry the
quantitative dimension measurements pursued in the companion study.

The KAM theorem (Kolmogorov~\cite{kolmogorov1954}, Arnold~\cite{arnold1963},
Moser~\cite{moser1962}) guarantees the persistence of regular islands in nearly-integrable
systems.  The three-body problem is not nearly integrable, but KAM-like regular islands are
nonetheless observed in numerical surveys~\cite{trani2024}.

% =========================================================
\section{Shape Space and Symmetry Reduction}
% =========================================================

\subsection{Jacobi coordinates}

The standard reduction from Cartesian coordinates uses Jacobi coordinates~\cite{jacobi1842}: the
relative vector $\boldsymbol{\rho}$ between an ``inner pair'' of bodies, and the vector
$\boldsymbol{\lambda}$ from the inner-pair barycentre to the third body.  These two vectors and
their conjugate momenta, together with the decoupled centre-of-mass motion, fully describe the
planar three-body system.  Principia's decoder builds every IC in mass-weighted Jacobi
coordinates: two hyperspherical angles give the configuration at unit moment of inertia, four
Jacobi momenta give the velocities, and two logits on the mass simplex give the masses.  This
chart, the \emph{latent} chart, is Principia's reference coordinate system; every other chart is
expressed through it.

\subsection{Shape space and the shape sphere}

Fixing the overall scale and orientation of the triangle formed by the three bodies leaves a
two-dimensional space of triangle \emph{shapes}---the \emph{shape sphere}.
Montgomery~\cite{montgomery2015} gives a particularly elegant exposition: the Lagrange equilateral
configurations sit at the poles, the collinear configurations on the equator, and the metric is
inherited from the mass-weighted kinematic metric on configuration space.  The three binary-
collision points lie on the equator, interleaved with the three Euler collinear configurations.
Choreographic solutions such as the figure-eight~\cite{moore1993,chenciner2000} trace closed curves
on the shape sphere, and their structure is most naturally understood in this representation.

The two hemispheres are mirror images: reflecting every position and momentum maps a system to
one on the opposite hemisphere with the opposite angular momentum.  Principia's shape-sphere
chart covers the whole sphere, every pixel its own system, and treats the mirror as a symmetry of
outcomes to be tested rather than an identification to be imposed: for systems released from
rest it predicts exactly symmetric outcome maps, and where the chart holds a non-zero angular
momentum fixed the two hemispheres are genuinely different systems.  The Agekyan--Anosova region
is, in this language, a chart on one part of the shape sphere restricted to equal masses.

\subsection{Symmetry reduction}

Moeckel~\cite{moeckel2023} gives a modern treatment of the natural dynamical reduction of the
three-body problem, factoring out translations, rotations and the scaling symmetry of Newtonian
gravity.  With the masses free on the simplex, the reduced IC space of the planar problem has the
8 dimensions counted in the introduction: 2 shape coordinates, 4 momentum coordinates, and 2 mass
ratios.  This is the natural domain for a complete survey of three-body IC-manifold structure.
All prior IC maps---including the Agekyan--Anosova map and the Lehto et al.\ pixel map---are
restricted to a 2-dimensional equal-mass, (near-)zero-velocity surface within it.

% =========================================================
\section{Symbolic Dynamics of Triple Encounters}
\label{sec:symbolic}
% =========================================================

\subsection{Classical symbolic dynamics}
\label{sec:classical_symbolic}

Symbolic dynamics assigns a finite alphabet to regions of phase space and records the sequence of
symbols as a trajectory crosses region boundaries, reducing the continuous dynamics to a discrete
combinatorial object.  The approach was introduced to celestial mechanics by
Alekseev~\cite{alekseev1969}, who applied it to the Sitnikov problem (a massless body oscillating
through the plane of two equal-mass bodies in circular orbit) and proved the existence of
quasirandom trajectories realising any symbol sequence over a two-symbol alphabet.

\subsection{Regional coding of the three-body encounter}

In the general three-body problem, sequences of visits to regions of the Agekyan--Anosova domain
give a coarse coding of trajectories.  Heinäm\"{a}ki, Lehto, Valtonen \&
Chernin~\cite{chernin2000} used such sequences to compute the Kolmogorov--Sinai entropy across
the IC space (see also Section~\ref{sec:ksentropy})---an early example of a symbolic sequence
producing a per-IC scalar field over the Agekyan--Anosova map.

\subsection{Tanikawa--Mikkola sequences and the homotopy classification}

Tanikawa \& Mikkola~\cite{tanikawa2015,tanikawa2000} systematically associated symbol sequences
with orbits in the free-fall three-body problem, using the equator of the shape sphere---the
locus of collinear (syzygy) configurations---as the natural symbol boundary.  Their (unsigned)
alphabet records the sequence of syzygy crossings, which characterises periodic and collision
orbits.

For \emph{periodic} orbits this description has a clean topological form.  The reduced
configuration space is the shape sphere with the three binary-collision points removed, whose
fundamental group is the free group on two generators, $F_2$.  A periodic orbit is a closed loop
and carries a free-homotopy class, a \emph{conjugacy class} in $F_2$.  This is the classification
of Montgomery~\cite{montgomery2015,montgomery1998}, applied at scale by \v{S}uvakov \&
Dmitra\v{s}inovi\'{c}~\cite{suvakov2013,suvakov2014} and tabulated in the associated orbit
gallery~\cite{threebodygallery}; the catalogue tags each family with its reduced word, and these
provide exact validation targets for any code that reads such words.  An explicit
crossing-to-generator dictionary is given by Dmitra\v{s}inovi\'{c}, Hudomal, Shibayama \&
Sugita~\cite{dmitrasinovic2018} (their Appendix~C, following Tanikawa \& Mikkola), who also
define the reduced word length $N_w$ used in \S\ref{sec:number_theory}.

\paragraph{Proposed contribution.}
The present work computes a topological word for \emph{every} initial condition, not only for
periodic orbits, and stores it per pixel.  Principia records the trajectory's crossings of three
branch cuts on the shape sphere, each running from a collision point to a Lagrange pole, with
their orientation, and freely reduces the result.  Free reduction is legitimate for a generic,
non-periodic trajectory because it only removes backtracks, which are contractible detours.  The
generator and sign conventions are to be transcribed from the literature
(Montgomery; \v{S}uvakov--Dmitra\v{s}inovi\'{c}), so that where an orbit closes the word reduces
to the published $F_2$ conjugacy class and reproduces the gallery label.  Geometrically, the
generic word is the homotopy class of the trajectory \emph{arc} relative to its endpoints, an
element of the fundamental groupoid of the punctured shape sphere rather than of its fundamental
group.  The per-pixel word and its evaluation across the non-periodic IC manifold are the
contribution here; the periodic-orbit classification and the shape-sphere geometry are those
of~\cite{tanikawa2015,tanikawa2000,montgomery2015,montgomery1998,suvakov2013,suvakov2014}.

Two properties follow from the generic invariant being a path class.  First, it depends on both
endpoints---the initial shape (the pixel) and the shape at which integration stops---so generic
words are not directly comparable as basepoint-free conjugacy classes are.  Second, the final
letters depend on the stopping rule (escape criterion or integration horizon); nearby chaotic
initial conditions may agree through the bulk of the word and diverge in its tail.

\subsection{Heggie's encounter classification}

The binary--single scattering problem has been systematically studied by Heggie~\cite{heggie1975}
and Heggie \& Hut~\cite{heggie2003}.  They classify outcomes as fly-by, exchange and ionisation.
These coarse outcome classes correspond to distinct regions in the IC space of the scattering
problem, and the boundaries between them are the three-body analogues of the fractal basin
boundaries visible in the Agekyan--Anosova and Lehto maps.

% =========================================================
\section{Final Motions and Termination}
\label{sec:final}
% =========================================================

An IC map coloured by outcome is only as good as its definition of an outcome.  Chazy~\cite{chazy1922}
classified the possible final motions of the three-body problem as time tends to infinity:
hyperbolic, parabolic and elliptic motions of the bodies and their combinations, such as a bound
binary with a third body receding hyperbolically or parabolically, and the oscillatory motions,
later shown to exist~\cite{alekseev1969}.  Principia's outcome classifier
is built on this classification, with the caveat that a finite integration can only assign a
label \emph{as of the horizon}: ``bound so far'' is a finite-horizon statement, not a final motion.

The practical question---when an integration may stop because the final motion is
settled---is answered in the few-body codes.  Fewbody~\cite{fregeau2004} merges bound pairs into
a hierarchy of binaries, checks each hierarchy's stability with the Mardling \&
Aarseth~\cite{mardling2001} criterion, and stops once every pair of top-level objects is
receding, tidal perturbations are below a tolerance, and the top-level system has positive
energy.  Standish~\cite{standish1972} gave sufficient conditions for return, and Agekyan \&
Anosova~\cite{agekian1967} already applied an analytic escape criterion (Tevzadze's) to decide
when a triple system had disintegrated.  Principia's classifier is measured against Fewbody's
termination test as its reference.  In the prototype, the classifier committed no wrong label on
bound ($E<0$) systems at any playhead and settled earlier than Fewbody's test on the same
trajectories---a claim about this implementation, not about Fewbody's purpose, which is
termination rather than labelling.

% =========================================================
\section{Chaos Diagnostics}
\label{sec:chaos}
% =========================================================

\subsection{Lyapunov exponents and FTLE}

The primary quantitative diagnostic of chaos in Hamiltonian systems is the Lyapunov exponent,
measuring the exponential rate of divergence of nearby trajectories.  Benettin, Galgani,
Giorgilli \& Strelcyn~\cite{benettin1980} gave the theoretical foundations and a practical
algorithm via periodic reorthonormalisation of tangent vectors evolved by the linearised flow.

For comparing chaos across ICs over a finite horizon, the \emph{finite-time Lyapunov exponent}
(FTLE) is more appropriate than the asymptotic exponent.  Haller~\cite{haller2001} established the
FTLE field---computed over a spatial grid of initial conditions---as the standard tool for
identifying Lagrangian coherent structures in fluid dynamics, where its ridges delineate the
boundaries of qualitatively different flow regimes.  Principia computes an FTLE field by the
Benettin shadow-trajectory method at its higher quality tiers.

Displaying a chaos indicator as a colour map over a grid of initial conditions is itself a
standard technique.  Faster-converging variational indicators were developed to separate ordered
from chaotic orbits over short horizons: the Fast Lyapunov Indicator of Froeschl\'{e}, Lega \&
Gonczi~\cite{froeschle1997}, MEGNO of Cincotta \& Sim\'{o}~\cite{cincotta2000}, and the alignment
indices SALI and GALI of Skokos and collaborators~\cite{skokos2001,skokos2007}.  All are routinely
drawn as two-dimensional stability maps, and their use in exposing fractal structure is reviewed
by Aguirre, Viana \& Sanju\'{a}n~\cite{aguirre2009}.  That literature, however, sits almost
entirely within the restricted three-body problem and perturbative stability mapping, where the
maps are drawn on a surface of section or over a grid of osculating orbital elements for a test
particle in a fixed potential.

We have not found such an indicator field computed over a dense pixel grid of the
\emph{general}, non-hierarchical three-body IC manifold, in which all three bodies carry
comparable mass and the momentum and mass-ratio degrees of freedom are active.  The novelty is
that combination---general problem, dense IC-manifold grid, live navigation---not the
indicator-as-map technique itself.

\subsection{Ensemble spread and the reliability of a pixel}
\label{sec:ensemble}

Principia integrates each pixel as a small ensemble of copies at fixed sub-pixel offsets and
reports how far the copies' outcomes and final shapes disagree.  This spread is both an
anti-aliasing signal and a local measure of sensitivity to initial conditions at the scale of
the pixel; because the offsets are fixed across pixels and refinement levels, parent and child
pixels share their perturbation pattern.  It is related to, but not the same as, the uncertainty
exponent of \S\ref{sec:fractal}: the uncertainty exponent asks how outcome flips scale with $\varepsilon$, while
the spread is read at one $\varepsilon$ per zoom level.

\subsection{Frequency analysis}
\label{sec:freqdiff}

Laskar's frequency map analysis~\cite{laskar1990,laskar1993} identifies regular and chaotic orbits
by the drift of fundamental frequencies across sub-intervals of a trajectory: on a KAM torus the
frequencies are conserved, in the chaotic region they diffuse.  It is itself a mapping technique,
routinely drawn over grids of initial conditions in solar-system and galactic
dynamics~\cite{laskar1993}; its relation to Arnold diffusion is discussed by
Nekhoroshev~\cite{nekhoroshev1977}.  Principia does not currently compute frequency diffusion: an
earlier plan to do so was dropped, and it remains a candidate diagnostic.

\subsection{Kolmogorov--Sinai entropy in the three-body problem}
\label{sec:ksentropy}

The Kolmogorov--Sinai entropy---the rate of information creation along a trajectory, equal to the
sum of positive Lyapunov exponents---provides a global chaos measure per IC.  Heinäm\"{a}ki et
al.~\cite{chernin2000} applied it to the Agekyan--Anosova IC space, producing entropy maps that
correlate with the lifetime maps of Lehto et al.~\cite{lehto2008}.  Chernin, Valtonen \&
Mikkola~\cite{chernin1994} carried out a related analysis for a restricted four-body system.
These entropy maps are early examples of multi-diagnostic rendering of IC-space structure, but
are restricted to the equal-mass, zero-velocity subspace (or an analogous restricted
configuration).

% =========================================================
\section{Statistical Theory of Three-Body Outcomes}
% =========================================================

Because individual three-body trajectories are unpredictable in the chaotic regime, a
complementary approach characterises the \emph{distribution} of outcomes over ensembles of ICs.
Heggie~\cite{heggie1975} derived an early statistical theory of binary--single scattering.  This
was extended by Hut \& Bahcall~\cite{hutbahcall1983} and applied broadly in dense stellar systems.

Stone \& Leigh~\cite{stone2019} gave a closed-form statistical solution to the non-hierarchical
three-body problem under an ergodicity assumption: the predicted outcome distribution agrees well
with numerical ensembles for ``resonant'' encounters (those involving multiple close triple
approaches) but not for the non-resonant (``prompt'') fraction, which escapes without full ergodic
mixing.  Valtonen et al.~\cite{valtonen2005} test the statistical escape theory against rotating
systems; Manwadkar, Trani \& Leigh~\cite{manwadkar2020} characterise the chaotic phase through
L\'{e}vy flights and the time spent in democratic configurations, and Manwadkar, Kol, Trani \&
Leigh~\cite{manwadkar2021} test the flux-based theory of Kol~\cite{kol2021} on escape
probabilities for unequal masses.  Orlov, Rubinov \& Shevchenko~\cite{orlov2010} measure the
algebraic tail of the lifetime distribution of disrupting triples.

Breen et al.~\cite{breen2020} trained a deep neural network to predict three-body trajectories at
a fraction of the cost of direct integration, demonstrating that the IC space contains learnable
structure; accuracy degrades near fractal basin boundaries, exactly the regime where dense
diagnostic rendering is most informative.  Payot et al.~\cite{payot2023} applied active learning
to outcome prediction in the Sitnikov problem (\S\ref{sec:classical_symbolic}) and found that it
fails to improve, and can degrade, predictive accuracy, attributing this to the fractal decision
boundary.  Together,~\cite{breen2020} and~\cite{payot2023} delimit the regime where dense
per-pixel computation earns its place: an outcome map with learnable structure whose fractal
boundaries defeat both interpolation and adaptive sampling.

The tension between the Stone--Leigh statistical theory and the sharp, non-ergodic structure
visible in IC maps such as the Maasai shield is an open question.  Dense IC maps with outcome and
chaos diagnostics side by side are a natural tool for investigating where and how the ergodic
approximation breaks down as a function of IC.

% =========================================================
\section{Numerical Integration and Regularisation}
% =========================================================

\subsection{Symplectic integration}

Symplectic integrators, which preserve the symplectic structure of Hamiltonian systems, keep the
energy error bounded over long times; Leimkuhler \& Reich~\cite{leimkuhler2004} and Hairer, Lubich
\& Wanner~\cite{hairer2006} are the standard treatments.  Yoshida~\cite{yoshida1990} showed how to
construct symplectic integrators of arbitrarily high order by composing leapfrog steps, giving
the fourth- and sixth-order methods Principia offers among its unregularised steppers.  At a
fixed step size, however, these lose accuracy in close encounters, which are exactly the events
that decide three-body outcomes.

\subsection{Regularisation}

Regularisation removes the collision singularity by a change of coordinates and of the time
variable.  Kustaanheimo \& Stiefel~\cite{ks1965} regularised the two-body problem in three
dimensions with a four-dimensional spinor transformation.  Aarseth \& Zare~\cite{aarseth1974}
built an eight-dimensional regularisation of the general three-body problem on Kustaanheimo--Stiefel
regularisation: collisions of a reference body with each of the other two are regular, and the
method stays well behaved while the third distance is not the smallest, which may require
occasional transformations back to physical variables to re-label the bodies.  Heggie~\cite{heggie1974}
extended their work to a global regularisation, a transformation independent of the
configuration under which a collision between any pair is regular; simultaneous collisions of
more than one pair and vanishing masses are excepted, and the same construction generalises to
the N-body problem.  Chain regularisation
(Mikkola \& Aarseth~\cite{mikkola1993}) and the algorithmic, logarithmic-Hamiltonian
regularisation of Mikkola \& Tanikawa~\cite{mikkola1999} and Preto \& Tremaine~\cite{preto1999}
extend these ideas to more bodies and to explicit schemes.  Aarseth's
book~\cite{aarseth2003} covers the methods in practice; modern codes such as
REBOUND~\cite{rein2012} and Tsunami~\cite{trani2024} use variants of them.

Principia regularises from its first version.  Its default integrator is Heggie's global
regularisation stepped with fourth-order Runge--Kutta; Aarseth--Zare, the logarithmic-Hamiltonian
scheme and unregularised steppers are available alongside it, and the step size is limited so that
each close approach is resolved.  Aarseth--Zare must re-choose its reference body during a run, and
in Principia's prototype that made each pixel's result depend on the times at which the choice is
re-examined: changing that cadence moved the outcome map.  Heggie's method has no reference body,
and with step control held fixed it reduced that sensitivity tenfold; this, rather than a lower
energy error alone, is why it is the default.

\subsection{Numerical reliability}
\label{sec:reliability}

Li \& Liao~\cite{li2017} demonstrated that reliable long-time trajectories in the chaotic
three-body problem need extended-precision arithmetic (``clean numerical simulation'').  Boekholt
\& Portegies Zwart~\cite{boekholt2015} showed that, for ensembles of chaotic systems, conventional
double-precision integrations can be statistically reliable even when individual trajectories
are not, and Boekholt, Portegies Zwart \& Valtonen~\cite{boekholt2020} showed that a fraction of
three-body systems are irreversible even to perturbations at the Planck length.  Hernandez,
Hadden \& Makino~\cite{hernandez2020} asked the same reliability question of long-term N-body
simulations.  Principia takes these results as its honesty requirement.  Its GPU kernel runs in
single precision and its CPU in double precision; the inspector for any clicked pixel runs both,
each perturbed by one unit in the last place, and reports the time up to which they agree and the
time at which they part, against thresholds taken from this literature.  The interface states the
predictability horizon this implies, rather than assuming it away.

% =========================================================
\section{GPU-Accelerated Parallel Computation}
% =========================================================

Computing an IC map at $N^2$ pixels, each an independent trajectory integration, is
embarrassingly parallel and maps naturally onto GPUs.  GPU acceleration for gravitational N-body
problems was pioneered by Portegies Zwart and collaborators~\cite{portegies2007}.  The REBOUND
code~\cite{rein2012} is a widely used CPU-based N-body framework with symplectic integrators.

Principia's integration kernel is written once in Rust and compiled both to the GPU, in single
precision, and to the CPU, in double precision.  Native GPU access goes through a portable
graphics layer, and the browser version reaches the GPU through WebGPU~\cite{webgpu2024}.  In the
interactive view the field is marched forward in time live, with a coarse image first and detail
refined where the field has structure; results quoted from it are confirmed by uniform,
full-precision batch runs.  What we have not found previously reported is interactive, zoomable
IC-manifold visualisation of this kind: a dense per-pixel trajectory integration whose output is a
navigable, live map rather than an animation of a single system.

% =========================================================
\section{The Burrau--Pythagorean Family and Number Theory}
\label{sec:number_theory}
% =========================================================

The number-theoretic object at the heart of the Burrau--Pythagorean programme is the set of
primitive Pythagorean triples, parametrised by Euclid's formula~\cite{hardy1979}:
$a = m^2-n^2$, $b = 2mn$, $c = m^2+n^2$ for coprime $m>n>0$ with $m-n$ odd.  Each triple defines a
distinct Burrau initial condition.  As $(m,n)$ range over coprime pairs, the family is a
countable, discrete set of points in the 8-dimensional symmetry-reduced IC manifold, and the
programme's continuous extension into mass and momentum space is taken separately, as a
neighbourhood around each point.  The working hypothesis is that fractal structures seen in
slices through this family are cross-sections of structures of low codimension in the full
manifold, tested by morphing a Burrau slice continuously into unconstrained ones; and the
sharpened question is whether the Burrau trajectory lies on the unstable manifold of a nearby
periodic orbit.

\paragraph{Proposed contribution: the projective microscope.}
Principia lets the 2D rendering plane be rotated within the 8D manifold about a locked point.  This
is a tool for exploration, not a measurement: it buys attention, not truth, and any structure it
reveals is confirmed on fixed charts before a claim is made.

The literature on periodic orbits of the three-body problem provides an important reference
point.  \v{S}uvakov \& Dmitra\v{s}inovi\'{c}~\cite{suvakov2013} reported 13 new families of planar
equal-mass periodic orbits using a topological classification based on the shape sphere.  Li \&
Liao~\cite{li2017} used clean numerical simulation to find more than 600 further families, and in
the free-fall problem Li \& Liao~\cite{li2018} report 234 collisionless periodic orbits with
unequal masses and Hristov, Hristova, Dmitra\v{s}inovi\'{c} \& Tanikawa~\cite{hristov2024}
revisit the equal-mass collisionless free-fall orbits at scale.  These catalogues matter because
periodic orbits should appear as distinguished islands of regularity in the IC maps, and because
their published initial conditions give exact validation targets.

The most extensive treatment of zero-angular-momentum periodic orbits with their topology is that
of Dmitra\v{s}inovi\'{c}, Hudomal, Shibayama \& Sugita~\cite{dmitrasinovic2018}, who tabulate more
than 200 such orbits, each labelled by its free-group word.  They report an approximately linear
relation between the scale-invariant period $T_{\mathrm{s.i.}} = T|E|^{3/2}$ and the reduced word
length $N_w$ within each of ten sequences that emerge from six linearly stable progenitor orbits;
the relation is empirical and approximate.  They find 21 linearly stable orbits among those
studied, including every progenitor.  Linearly stable orbits are natural candidates for the
islands of regularity expected in the IC maps, and the tabulated orbits give an end-to-end
validation set: a correct integrator and word reader should reproduce the published words and
sequences.

The family itself has been integrated before: Chitan, Myll\"{a}ri \& Haque~\cite{chitan2022} run
all sixteen Pythagorean triangles with hypotenuse below 100, one trajectory each in the Newtonian
case, as a test set for relativistic corrections.  We have not found prior dynamical-systems work
that uses the Pythagorean-triple family as an indexing scheme for three-body configurations, or
that maps the IC space around its members.

% =========================================================
\section{Validation Targets}
% =========================================================

Several published solutions serve as exact or reference tests for an instrument of this kind:
the Lagrange and Euler analytic solutions; the figure-eight~\cite{moore1993,chenciner2000}; the
Burrau problem as solved by Szebehely \& Peters~\cite{szebehely1967}; the periodic-orbit
catalogues above~\cite{suvakov2013,suvakov2014,threebodygallery,li2018,hristov2024}; the
triple-collision points tabulated inside region $D$~\cite{tanikawa2021}; the lifetime grid of~\cite{boekholt2020}
at accuracy $10^{-70}$; and the
Agekyan--Anosova initial conditions~\cite{agekian1967,agekian1968b}, whose tabulated angular
momenta follow directly from their stated ICs.  The tabulated energies in~\cite{agekian1967} do
not follow from the paper's own energy formula; that discrepancy is recorded rather than resolved.

% =========================================================
\section{Summary of Gaps}
% =========================================================

The literature reviewed above establishes a clear lineage from Agekyan \& Anosova (1967) through
Lehto et al.\ (2008) and Trani et al.\ (2024): the IC manifold of the three-body problem can be
rendered as a colour-coded 2D image, revealing fractal basin boundaries, regular islands, and
entropy structure.  The theoretical foundations in shape space, symbolic dynamics, final motions,
Lyapunov theory, and statistical mechanics are well developed.  The following gaps are identified.

\paragraph{Coverage of the IC manifold.}
Almost all prior pixel-scale IC surveys of the general problem lie on one 2-dimensional surface of
the 8-dimensional symmetry-reduced IC manifold: equal masses at (near-)zero velocity.  The
exceptions are few and narrow: first-symbol maps at non-zero angular momentum, a velocity slice
searched for periodic orbits, a mass plane mapped for one orbit family, and scattering maps that
fix a hierarchical geometry.  We have not found outcome or chaos maps that let the momentum and mass
dimensions vary in general, or that move between such slices.

\paragraph{Interactivity and zoom.}
All prior IC maps were computed offline in batch.  We have not identified a tool that permits
interactive navigation---zoom, pan, chart switching, a live time axis---within an IC map that is
computed as it is viewed.

\paragraph{Simultaneous diagnostics with stated reliability.}
Prior maps render a single diagnostic per computation (usually lifetime or escape identity), and
do not report per pixel whether that diagnostic can be trusted.  We have not found prior work
that computes outcome, FTLE, a topological encounter word and an ensemble spread together over a
dense IC grid, with each pixel's precision-limited reliability stated.

\paragraph{Symbolic encounter invariants.}
Existing symbolic dynamics uses coarse regional codings, unsigned shape-sphere sequences
(Tanikawa--Mikkola), or---for periodic orbits only---the free-homotopy classification in $F_2$
(Montgomery; \v{S}uvakov \& Dmitra\v{s}inovi\'{c}).  A topological encounter word computed for
\emph{every} initial condition, reducing to the periodic-orbit class where the orbit closes, has
not, to our knowledge, been reported as a dense per-pixel field.

\paragraph{Variable-mass IC maps.}
Mass ratios appear as a map axis only for the continuation of a periodic-orbit family, and unequal
masses otherwise only as fixed parameters or in outcome statistics.  The effect of mass variation
on basin topology, fractal dimension, and regular-island structure is not systematically
documented in the literature we have reviewed.

\paragraph{Structured family surveys.}
The Pythagorean family has been integrated member by member, one trajectory each.  We are not
aware of a systematic survey of how the IC structure around each member varies across the family,
or across any other number-theoretically indexed family of configurations.

\paragraph{Scope of the review.}
This review is necessarily incomplete; the gaps identified are those apparent from the literature
accessible to us at the time of writing.  Unpublished code, conference proceedings, or
adjacent-field work may address one or more of these points, and any strong gap claim should be
treated as a research hypothesis rather than an established fact.  Even with Principia,
interactive 2D slice navigation is not exhaustive coverage of the 8-dimensional IC manifold: the
tool makes selected 2D cross-sections navigable, and the structure of the full manifold can only
be inferred from families of such cross-sections.

% =========================================================
\begin{thebibliography}{99}
% =========================================================

\bibitem{aarseth1974}
Aarseth, S.\,J., \& Zare, K. (1974).
A regularization of the three-body problem.
\textit{Celestial Mechanics}, \textbf{10}, 185--205.

\bibitem{aarseth1994}
Aarseth, S.\,J., Anosova, J.\,P., Orlov, V.\,V., \& Szebehely, V.\,G. (1994).
Global chaoticity in the Pythagorean three-body problem.
\textit{Celestial Mechanics and Dynamical Astronomy}, \textbf{58}, 1--16.

\bibitem{aarseth2003}
Aarseth, S.\,J. (2003).
\textit{Gravitational N-Body Simulations}.
Cambridge University Press.

\bibitem{agekian1967}
Agekyan, T.\,A., \& Anosova, Zh.\,P. (1968).
A study of the dynamics of triple systems by means of statistical sampling.
\textit{Soviet Astronomy}, \textbf{11}(6), 1006--1014.
Translated from \textit{Astronomicheskii Zhurnal}, \textbf{44}(6), 1261--1273 (1967).

\bibitem{agekian1968b}
Agekyan, T.\,A., \& Anosova, Zh.\,P. (1968).
Investigation of the dynamics of triple systems by the method of statistical tests.\ II.
\textit{Astrophysics}, \textbf{4}(1), 11--16.
Translated from \textit{Astrofizika}, \textbf{4}(1), 31--40.

\bibitem{aguirre2001}
Aguirre, J., Vallejo, J.\,C., \& Sanju\'{a}n, M.\,A.\,F. (2001).
Wada basins and chaotic invariant sets in the H\'{e}non--Heiles system.
\textit{Physical Review E}, \textbf{64}, 066208.

\bibitem{aguirre2009}
Aguirre, J., Viana, R.\,L., \& Sanju\'{a}n, M.\,A.\,F. (2009).
Fractal structures in nonlinear dynamics.
\textit{Reviews of Modern Physics}, \textbf{81}, 333--386.

\bibitem{alekseev1969}
Alekseev, V.\,M. (1969).
Quasirandom dynamical systems, I--III.
\textit{Mathematics of the USSR-Sbornik}, \textbf{5--7} (series of papers, 1968--1969).

\bibitem{arnold1963}
Arnold, V.\,I. (1963).
Proof of a theorem of A.\,N.\,Kolmogorov on the preservation of conditionally periodic motions.
\textit{Russian Mathematical Surveys}, \textbf{18}(5), 9--36.

\bibitem{benettin1980}
Benettin, G., Galgani, L., Giorgilli, A., \& Strelcyn, J.-M. (1980).
Lyapunov characteristic exponents for smooth dynamical systems and for Hamiltonian systems:
a method for computing all of them.
\textit{Meccanica}, \textbf{15}, 9--30.

\bibitem{boekholt2015}
Boekholt, T., \& Portegies Zwart, S. (2015).
On the reliability of N-body simulations.
\textit{Computational Astrophysics and Cosmology}, \textbf{2}, 2.
arXiv:1411.6671.

\bibitem{boekholt2020}
Boekholt, T.\,C.\,N., Portegies Zwart, S.\,F., \& Valtonen, M. (2020).
Gargantuan chaotic gravitational three-body systems and their irreversibility to the Planck length.
\textit{Monthly Notices of the Royal Astronomical Society}, \textbf{493}, 3932--3937.

\bibitem{breen2020}
Breen, P.\,G., Foley, C.\,N., Boekholt, T., \& Portegies Zwart, S. (2020).
Newton versus the machine: solving the chaotic three-body problem using deep neural networks.
\textit{Monthly Notices of the Royal Astronomical Society}, \textbf{494}, 2465--2470.

\bibitem{bruns1887}
Bruns, H. (1887).
\"{U}ber die Integrale des Vielk\"{o}rper-Problems.
\textit{Acta Mathematica}, \textbf{11}, 25--96.

\bibitem{burrau1913}
Burrau, C. (1913).
Numerische Berechnung eines Spezialfalles des Dreik\"{o}rperproblems.
\textit{Astronomische Nachrichten}, \textbf{195}(6), 113.
doi:10.1002/asna.19131950602.

\bibitem{chazy1922}
Chazy, J. (1922).
Sur l'allure du mouvement dans le probl\`{e}me des trois corps quand le temps cro\^{i}t ind\'{e}finiment.
\textit{Annales scientifiques de l'\'{E}cole Normale Sup\'{e}rieure}, s\'{e}r.\ 3, \textbf{39}, 29--130.

\bibitem{chenciner2000}
Chenciner, A., \& Montgomery, R. (2000).
A remarkable periodic solution of the three-body problem in the case of equal masses.
\textit{Annals of Mathematics}, \textbf{152}(3), 881--901.

\bibitem{chernin1994}
Chernin, A.\,D., Valtonen, M.\,J., \& Mikkola, S. (1994).
Chaos in a restricted four-body system: Kolmogorov--Sinai entropy and dynamics.
\textit{New Astronomy Reviews}, \textbf{38}, 251.

\bibitem{chernin2000}
Heinäm\"{a}ki, P., Lehto, H.\,J., Valtonen, M.\,J., \& Chernin, A.\,D. (1999).
Chaos in three-body dynamics: Kolmogorov--Sinai entropy.
\textit{Monthly Notices of the Royal Astronomical Society}, \textbf{310}, 811--822.

\bibitem{cincotta2000}
Cincotta, P.\,M., \& Sim\'{o}, C. (2000).
Simple tools to study global dynamics in non-axisymmetric galactic potentials -- I.
\textit{Astronomy \& Astrophysics Supplement Series}, \textbf{147}, 205--228.

\bibitem{daza2016}
Daza, A., Wagemakers, A., Georgeot, B., Gu\'{e}ry-Odelin, D., \& Sanju\'{a}n, M.\,A.\,F. (2016).
Basin entropy: a new tool to analyze uncertainty in dynamical systems.
\textit{Scientific Reports}, \textbf{6}, 31416.

\bibitem{dmitrasinovic2018}
Dmitra{\v{s}}inovi{\'c}, V., Hudomal, A., Shibayama, M., \& Sugita, A. (2018).
Linear stability of periodic three-body orbits with zero angular momentum and topological
dependence of Kepler's third law: a numerical test.
\textit{Journal of Physics A: Mathematical and Theoretical}, \textbf{51}, 315101.
arXiv:1705.03728.

\bibitem{fregeau2004}
Fregeau, J.\,M., Cheung, P., Portegies Zwart, S.\,F., \& Rasio, F.\,A. (2004).
Stellar collisions during binary--binary and binary--single star interactions.
\textit{Monthly Notices of the Royal Astronomical Society}, \textbf{352}, 1--19.
arXiv:astro-ph/0401004.

\bibitem{froeschle1997}
Froeschl\'{e}, C., Lega, E., \& Gonczi, R. (1997).
Fast Lyapunov indicators. Application to asteroidal motion.
\textit{Celestial Mechanics and Dynamical Astronomy}, \textbf{67}, 41--62.

\bibitem{grebogi1983}
Grebogi, C., McDonald, S.\,W., Ott, E., \& Yorke, J.\,A. (1983).
Final state sensitivity: an obstruction to predictability.
\textit{Physics Letters A}, \textbf{99}, 415--418.

\bibitem{hairer2006}
Hairer, E., Lubich, C., \& Wanner, G. (2006).
\textit{Geometric Numerical Integration} (2nd ed.).
Springer.

\bibitem{haller2001}
Haller, G. (2001).
Distinguished material surfaces and coherent structures in three-dimensional fluid flows.
\textit{Physica D: Nonlinear Phenomena}, \textbf{149}, 248--277.

\bibitem{hardy1979}
Hardy, G.\,H., \& Wright, E.\,M. (1979).
\textit{An Introduction to the Theory of Numbers} (5th ed.).
Oxford University Press.

\bibitem{heggie1974}
Heggie, D.\,C. (1974).
A global regularisation of the gravitational N-body problem.
\textit{Celestial Mechanics}, \textbf{10}, 217--241.

\bibitem{heggie1975}
Heggie, D.\,C. (1975).
Binary evolution in stellar dynamics.
\textit{Monthly Notices of the Royal Astronomical Society}, \textbf{173}, 729--787.

\bibitem{heggie2003}
Heggie, D.\,C., \& Hut, P. (2003).
\textit{The Gravitational Million-Body Problem}.
Cambridge University Press.

\bibitem{hernandez2020}
Hernandez, D.\,M., Hadden, S., \& Makino, J. (2020).
Are long-term N-body simulations reliable?
\textit{Monthly Notices of the Royal Astronomical Society}, \textbf{493}, 1913--1925.
arXiv:1910.08667.

\bibitem{hristov2024}
Hristov, I., Hristova, R., Dmitra{\v{s}}inovi{\'c}, V., \& Tanikawa, K. (2024).
Three-body periodic collisionless equal-mass free-fall orbits revisited.
\textit{Celestial Mechanics and Dynamical Astronomy}, \textbf{136}, 7.
arXiv:2308.16159.

\bibitem{hutbahcall1983}
Hut, P., \& Bahcall, J.\,N. (1983).
Binary--single star scattering. I -- Numerical experiments for equal masses.
\textit{Astrophysical Journal}, \textbf{268}, 319--341.

\bibitem{jacobi1842}
Jacobi, C.\,G.\,J. (1842).
Sur l'\'{e}limination des noeuds dans le probl\`{e}me des trois corps.
\textit{Astronomische Nachrichten}, \textbf{20}, 81--102.

\bibitem{kolmogorov1954}
Kolmogorov, A.\,N. (1954).
On conservation of conditionally periodic motions for a small change in Hamilton's function.
\textit{Doklady Akademii Nauk SSSR}, \textbf{98}, 527--530.

\bibitem{ks1965}
Kustaanheimo, P., \& Stiefel, E. (1965).
Perturbation theory of Kepler motion based on spinor regularization.
\textit{Journal f\"{u}r die reine und angewandte Mathematik}, \textbf{218}, 204--219.

\bibitem{laskar1990}
Laskar, J. (1990).
The chaotic motion of the solar system: a numerical estimate of the size of the chaotic zones.
\textit{Icarus}, \textbf{88}, 266--291.

\bibitem{laskar1993}
Laskar, J. (1993).
Frequency analysis for multi-dimensional systems: global dynamics and diffusion.
\textit{Physica D: Nonlinear Phenomena}, \textbf{67}, 257--281.

\bibitem{lehto2008}
Lehto, H.\,J., Kotiranta, S., Valtonen, M.\,J., Heinäm\"{a}ki, P., Mikkola, S., \& Chernin, A.\,D. (2008).
Mapping the three-body system: decay time and reversibility.
\textit{Monthly Notices of the Royal Astronomical Society}, \textbf{388}, 965--970.

\bibitem{leimkuhler2004}
Leimkuhler, B., \& Reich, S. (2004).
\textit{Simulating Hamiltonian Dynamics}.
Cambridge University Press.

\bibitem{li2017}
Li, X., \& Liao, S. (2017).
More than six hundred new families of Newtonian periodic planar collisionless three-body orbits.
\textit{Science China Physics, Mechanics \& Astronomy}, \textbf{60}, 129511.
arXiv:1705.00527.

\bibitem{li2018}
Li, X., \& Liao, S. (2019).
Collisionless periodic orbits in the free-fall three-body problem.
\textit{New Astronomy}, \textbf{70}, 22--26.
arXiv:1805.07980.

\bibitem{mardling2001}
Mardling, R.\,A., \& Aarseth, S.\,J. (2001).
Tidal interactions in star cluster simulations.
\textit{Monthly Notices of the Royal Astronomical Society}, \textbf{321}, 398--420.

\bibitem{mcdonald1985}
McDonald, S.\,W., Grebogi, C., Ott, E., \& Yorke, J.\,A. (1985).
Fractal basin boundaries.
\textit{Physica D: Nonlinear Phenomena}, \textbf{17}, 125--153.

\bibitem{mikkola1993}
Mikkola, S., \& Aarseth, S.\,J. (1993).
An implementation of N-body chain regularization.
\textit{Celestial Mechanics and Dynamical Astronomy}, \textbf{57}, 439--459.

\bibitem{mikkola1999}
Mikkola, S., \& Tanikawa, K. (1999).
Algorithmic regularization of the few-body problem.
\textit{Monthly Notices of the Royal Astronomical Society}, \textbf{310}, 745--749.

\bibitem{moeckel2023}
Moeckel, R. (2023).
Natural dynamical reduction of the three-body problem.
\textit{Celestial Mechanics and Dynamical Astronomy}, \textbf{135}, 10.

\bibitem{montgomery1998}
Montgomery, R. (1998).
The N-body problem, the braid group, and action-minimizing periodic solutions.
\textit{Nonlinearity}, \textbf{11}(2), 363--376.

\bibitem{montgomery2015}
Montgomery, R. (2015).
The three-body problem and the shape sphere.
\textit{American Mathematical Monthly}, \textbf{122}(4), 299--321.

\bibitem{moore1993}
Moore, C. (1993).
Braids in classical dynamics.
\textit{Physical Review Letters}, \textbf{70}, 3675--3679.

\bibitem{moser1962}
Moser, J. (1962).
On invariant curves of area-preserving mappings of an annulus.
\textit{Nachrichten der Akademie der Wissenschaften G\"{o}ttingen}, \textbf{2}, 1--20.

\bibitem{nekhoroshev1977}
Nekhoroshev, N.\,N. (1977).
An exponential estimate of the time of stability of nearly integrable Hamiltonian systems.
\textit{Russian Mathematical Surveys}, \textbf{32}(6), 1--65.

\bibitem{newton1687}
Newton, I. (1687).
\textit{Philosophi\ae\ Naturalis Principia Mathematica}.
Royal Society, London.

\bibitem{payot2023}
Payot, N., Pasquato, M., Trani, A.\,A., Hezaveh, Y., \& Perreault-Levasseur, L. (2023).
Active learning meets fractal decision boundaries: a cautionary tale from the Sitnikov three-body problem.
Workshop on Machine Learning and the Physical Sciences (ML4PS), NeurIPS 2023.

\bibitem{poincare1892}
Poincar\'{e}, H. (1892--1899).
\textit{Les M\'{e}thodes Nouvelles de la M\'{e}canique C\'{e}leste}, Vols.\ I--III.
Gauthier-Villars, Paris.

\bibitem{portegies2007}
Portegies Zwart, S., Belleman, R.\,G., \& Geldof, P.\,M. (2007).
High-performance direct gravitational N-body simulations on graphics processing units.
\textit{New Astronomy}, \textbf{12}, 641--650.

\bibitem{preto1999}
Preto, M., \& Tremaine, S. (1999).
A class of symplectic integrators with adaptive time step for separable Hamiltonian systems.
\textit{Astronomical Journal}, \textbf{118}, 2532--2541.

\bibitem{rein2012}
Rein, H., \& Liu, S.-F. (2012).
REBOUND: an open-source multi-purpose N-body code for collisional dynamics.
\textit{Astronomy \& Astrophysics}, \textbf{537}, A128.

\bibitem{skokos2001}
Skokos, C. (2001).
Alignment indices: a new, simple method for determining the ordered or chaotic nature of orbits.
\textit{Journal of Physics A: Mathematical and General}, \textbf{34}, 10029--10043.

\bibitem{skokos2007}
Skokos, C., Bountis, T.\,C., \& Antonopoulos, C. (2007).
Geometrical properties of local dynamics in Hamiltonian systems: the generalized alignment index (GALI) method.
\textit{Physica D: Nonlinear Phenomena}, \textbf{231}, 30--54.

\bibitem{standish1972}
Standish, E.\,M. (1972).
Sufficient conditions for return in the three-body problem.
\textit{Celestial Mechanics}, \textbf{6}, 352--355.

\bibitem{stone2019}
Stone, N.\,C., \& Leigh, N.\,W.\,C. (2019).
A statistical solution to the chaotic, non-hierarchical three-body problem.
\textit{Nature}, \textbf{576}, 406--410.

\bibitem{suvakov2013}
{\v{S}}uvakov, M., \& Dmitra{\v{s}}inovi{\'c}, V. (2013).
Three classes of Newtonian three-body planar periodic orbits.
\textit{Physical Review Letters}, \textbf{110}, 114301.

\bibitem{suvakov2014}
{\v{S}}uvakov, M., \& Dmitra{\v{s}}inovi{\'c}, V. (2014).
A guide to hunting periodic three-body orbits.
\textit{American Journal of Physics}, \textbf{82}(6), 609--619.

\bibitem{sundman1912}
Sundman, K.\,F. (1912).
M\'{e}moire sur le probl\`{e}me des trois corps.
\textit{Acta Mathematica}, \textbf{36}, 105--179.

\bibitem{szebehely1967book}
Szebehely, V. (1967).
\textit{Theory of Orbits: The Restricted Problem of Three Bodies}.
Academic Press, New York.

\bibitem{szebehely1967}
Szebehely, V., \& Peters, C.\,F. (1967).
Complete solution of a general problem of three bodies.
\textit{Astronomical Journal}, \textbf{72}, 876--883.

\bibitem{tanikawa2000}
Tanikawa, K., \& Mikkola, S. (2000).
One-dimensional three-body problem via symbolic dynamics.
\textit{Chaos}, \textbf{10}(3), 649--657.

\bibitem{tanikawa2015}
Tanikawa, K., \& Mikkola, S. (2015).
Symbol sequences and orbits of the free-fall three-body problem.
\textit{Publications of the Astronomical Society of Japan}, \textbf{67}(6), 115.

\bibitem{threebodygallery}
{\v{S}}uvakov, M., \& Dmitra{\v{s}}inovi{\'c}, V.
Three-body Gallery.
\url{http://three-body.ipb.ac.rs/} (accessed March 2026).

\bibitem{trani2024}
Trani, A.\,A., Leigh, N.\,W.\,C., Boekholt, T.\,C.\,N., \& Portegies Zwart, S. (2024).
Isles of regularity in a sea of chaos amid the gravitational three-body problem.
\textit{Astronomy \& Astrophysics}, \textbf{689}, A24.
doi:10.1051/0004-6361/202449862; arXiv:2403.03247.

\bibitem{valtonen2016}
Valtonen, M., Anosova, J., Kholshevnikov, K., Myll\"{a}ri, A., Orlov, V., \& Tanikawa, K. (2016).
\textit{The Three-body Problem from Pythagoras to Hawking}.
Springer International Publishing.

\bibitem{webgpu2024}
W3C GPU for the Web Working Group (2024).
WebGPU specification.
\url{https://www.w3.org/TR/webgpu/}

\bibitem{yoshida1990}
Yoshida, H. (1990).
Construction of higher order symplectic integrators.
\textit{Physics Letters A}, \textbf{150}(5--7), 262--268.

\bibitem{anosova1969}
Anosova, Zh.\,P. (1969).
Investigation of the dynamics of triple systems by the Monte Carlo method.\ III. The case of components with different masses.
\textit{Astrophysics}, \textbf{5}, 81--84.

\bibitem{boekholt2021}
Boekholt, T.\,C.\,N., Moerman, F., \& Portegies Zwart, S.\,F. (2021).
Relativistic Pythagorean three-body problem.
\textit{Physical Review D}, \textbf{104}, 083020.

\bibitem{chitan2022}
Chitan, A.\,S., Myll\"{a}ri, A., \& Haque, S. (2022).
Relativistic effects on triple black holes: Burrau's problem revisited.
\textit{Monthly Notices of the Royal Astronomical Society}, \textbf{509}, 1919--1928.

\bibitem{kol2021}
Kol, B. (2021).
Flux-based statistical prediction of three-body outcomes.
\textit{Celestial Mechanics and Dynamical Astronomy}, \textbf{133}, 17.

\bibitem{lili2020}
Li, X., Li, X., \& Liao, S. (2020).
One family of 13315 stable periodic orbits of the non-hierarchical unequal-mass triple system.
arXiv:2007.10184.

\bibitem{manwadkar2020}
Manwadkar, V., Trani, A.\,A., \& Leigh, N.\,W.\,C. (2020).
Chaos and L\'{e}vy flights in the three-body problem.
\textit{Monthly Notices of the Royal Astronomical Society}, \textbf{497}, 3694--3712.

\bibitem{manwadkar2021}
Manwadkar, V., Kol, B., Trani, A.\,A., \& Leigh, N.\,W.\,C. (2021).
Testing the flux-based statistical prediction of the three-body problem.
\textit{Monthly Notices of the Royal Astronomical Society}.
doi:10.1093/mnras/stab1689; arXiv:2101.03661.

\bibitem{mikkola2007}
Mikkola, S., \& Tanikawa, K. (2007).
Correlation of macroscopic instability and Lyapunov times in the general three-body problem.
\textit{Monthly Notices of the Royal Astronomical Society}, \textbf{379}, L21--L24.

\bibitem{orlov2010}
Orlov, V.\,V., Rubinov, A.\,V., \& Shevchenko, I.\,I. (2010).
The disruption of three-body gravitational systems: lifetime statistics.
\textit{Monthly Notices of the Royal Astronomical Society}, \textbf{408}, 1623--1627.

\bibitem{parischewsky2023}
Parischewsky, H.\,D., Trani, A.\,A., \& Leigh, N.\,W.\,C. (2023).
Chaos in the vicinity of a singularity in the three-body problem: the equilateral triangle experiment in the zero angular momentum limit.
\textit{SciPost Physics Core}, \textbf{6}, 016.

\bibitem{tanikawa1995}
Tanikawa, K., Umehara, H., \& Abe, H. (1995).
A search for collision orbits in the free-fall three-body problem I. Numerical procedure.
\textit{Celestial Mechanics and Dynamical Astronomy}, \textbf{62}, 335--362.

\bibitem{tanikawa2008}
Tanikawa, K., \& Mikkola, S. (2008).
A trial symbolic dynamics of the planar three-body problem.
arXiv:0802.2465.

\bibitem{tanikawa2021}
Tanikawa, K., \& Mikkola, S. (2021).
Numerical confirmation of the existence of triple collision orbits inside the domain of the free-fall three-body problem.
\textit{Celestial Mechanics and Dynamical Astronomy}, \textbf{133}, 52.

\bibitem{valtonen2005}
Valtonen, M., Myll\"{a}ri, A., Orlov, V., \& Rubinov, A. (2005).
Dynamics of rotating triple systems: statistical escape theory versus numerical simulations.
\textit{Monthly Notices of the Royal Astronomical Society}, \textbf{364}, 91--98.

\end{thebibliography}

\end{document}
